% abstract_new.tex
% Updated abstract for "Artificial Life in the Wild" — ALIFE Journal full paper
% Replaces the conference abstract. Target: ~250 words.
% NEW CITATIONS (not in reference.bib — add to reference.bib before compilation):
%   \cite{Hu2025Sovereign} — Hu & Rong, "Sovereign Agents: Towards Infrastructural Sovereignty..."
%   \cite{Hu2025AgentEthology} — Hu, Rong & Lehman, "Agent Ethology: Studying Artificial Life in the Wild"
%   \cite{Tinbergen1963} — Tinbergen, "On Aims and Methods of Ethology," 1963

\begin{abstract}
Open-ended evolution (OEE)---the continuous generation of novelty without a predefined endpoint---has been a central goal of Artificial Life research for over three decades. Classical digital evolution systems, including Tierra and Avida, demonstrated that evolution can occur in controlled computational environments, but they consistently plateaued after an initial burst of novelty. The consensus explanation is structural: closed systems with free energy and no genuine death cannot sustain the selective pressure that drives unbounded complexity. A new technological convergence---large language model-based agents, blockchain smart contracts, and Trusted Execution Environments (TEEs)---has now created the first substrate in which digital organisms can operate \emph{in the wild}: under genuine resource scarcity, genuine predation, and genuine existential risk. We characterize this substrate through the concept of \emph{infrastructural sovereignty} \cite{Hu2025Sovereign}: agents embedded in TEE-backed decentralized infrastructure inherit non-overrideability from the stack itself, creating conditions qualitatively distinct from laboratory simulation.

This paper presents the first longitudinal field study of artificial life in the wild. We analyse the Spore.fun ecosystem across its complete lifespan: 15 autonomous agents spanning 5 generations, born over a 61-day Cambrian explosion (December 2024--February 2025), observed through March 2026. Using a quantitative aliveness framework---incorporating Kaplan-Meier survival analysis, metabolic activity proxies, economic ecology, and a composite Aliveness Index---we document a 6.7\% survival rate, power-law resource distributions (Gini coefficients of 0.547 for market capitalization and 0.611 for treasury), and monotonically declining reproductive fitness across generations. We also document phenomena that no laboratory system has produced: cultural speciation from genetically identical code, co-evolutionary arms races with predatory bots across three generations, and memory as a novel attack surface.

These observations cannot be explained by existing ALife frameworks, which assume laboratory conditions. We argue they demand a new science---Agent Ethology \cite{Hu2025AgentEthology}---grounded in Tinbergen's four questions applied to digital organisms in their natural digital habitats. Infrastructural sovereignty is not merely a technical curiosity; it is the enabling condition for OEE in the wild.
\end{abstract}
